\ifwithsols
\else
\clearpage
\fi

\subsection*{سؤال 2 امتحان 899381 سنة 2023}
\addcontentsline{toc}{subsection}{سؤال 2 امتحان 899381 سنة 2023}

\noindent
\makebox[\textwidth][c]{\includegraphics[width=0.9\paperwidth,keepaspectratio]{ ../../../bagrut_questions/basics/arrays_2023_899381_2.png }%
}%

\begin{boxHint}
يمكنك معرفة الفرق الثابت بين الأعداد في المصفوفة من خلال مقارنة الفرق بين أول عنصرين مع الفرق بين العنصرين الثاني والثالث. 
(الفرق الثابت هو الفرق الأصغر بينهما، أو إذا كانا متساويين فهما مساويان للفرق الثابت) \\
بعد معرفة الفرق الثابت، يمكنك المرور على المصفوفة لإيجاد العنصرين اللذين الفرقُ بينهما أكبر من الفرق الثابت، وبالتالي حساب الحد الناقص.
\end{boxHint}


\ifwithsols
\begin{boxSolution}[سؤال 2 امتحان 899381 سنة 2023]

\begin{boxCode}
\begin{english}
\begin{minted}{csharp}
public static int MissingNum(int[] arr)
{
    int diff1 = arr[1] - arr[0];
    int diff2 = arr[2] - arr[1];
    int constantDiff;
    if (diff1 <= diff2)
        constantDiff = diff1;
    else
        constantDiff = diff2;

    for (int i = 0; i < arr.Length - 1; i++)
    {
        int currentDiff = arr[i + 1] - arr[i];
        if (currentDiff != constantDiff)
            return arr[i] + constantDiff;
    }
    return -1;
}
\end{minted}
\end{english}
\end{boxCode}

\end{boxSolution}
\clearpage
\fi
